% =========================================================
%  frontmatter/learning-outcomes.tex
%  Capaian pembelajaran dan peta kompetensi
% =========================================================

\chapter*{Capaian Pembelajaran dan Peta Kompetensi}
\addcontentsline{toc}{chapter}{Capaian Pembelajaran dan Peta Kompetensi}
\thispagestyle{plain}

Bagian ini memuat capaian pembelajaran yang dibebankan pada mata kuliah
\NamaMK\ beserta peta keterkaitannya dengan bab-bab pada buku ini. Pembaca
yang ingin mengetahui alasan sebuah materi diajarkan, atau ingin memastikan
materi mana yang menopang suatu asesmen, dapat berpangkal dari tabel-tabel
berikut.

\section*{Capaian Pembelajaran Lulusan (CPL)}

\begin{table}[htbp]
\centering
\caption{Rumusan capaian pembelajaran lulusan yang dibebankan pada mata kuliah}
\label{tab:cpl-buku}
\begin{tabularx}{\textwidth}{@{}C{1.8cm} Y@{}}
\toprule
\headrow \textbf{Kode} & \textbf{Rumusan} \\
\midrule
\CPL{03} &
Kemampuan merancang dan menganalisis algoritma untuk menyelesaikan
permasalahan secara optimal, serta memilih dan menerapkannya pada bahasa
pemrograman tertentu. \\
\addlinespace
\CPL{04} &
Kemampuan mendesain, mengimplementasi, dan mengevaluasi solusi berbasis
\istilah{computing} yang memenuhi kebutuhan \istilah{computing} pada sebuah
disiplin program. \\
\bottomrule
\end{tabularx}
\end{table}

\section*{Capaian Pembelajaran Mata Kuliah (CPMK)}

\begin{table}[htbp]
\centering
\caption{Rumusan CPMK, CPL yang didukung, dan ranah taksonomi Bloom}
\label{tab:cpmk-buku}
\begin{tabularx}{\textwidth}{@{}C{1.7cm}Y C{3.4cm} C{1.7cm}@{}}
\toprule
\headrow \textbf{Kode} & \textbf{Rumusan CPMK} & \textbf{CPL Terkait} & \textbf{Bloom} \\
\midrule
\CPMK{1} & Merancang \istilah{game loop} dan logika fisika sederhana. &
\CPL{03} (utama), \CPL{04} (pendukung) & C2--C4 \\
\addlinespace
\CPMK{2} & Menerapkan algoritma kecerdasan buatan (AI) dasar untuk NPC. &
\CPL{03} (utama), \CPL{04} (pendukung) & C3--C5 \\
\addlinespace
\CPMK{3} & Menggunakan \istilah{game engine} (Unity) untuk membangun level. &
\CPL{04} (utama), \CPL{03} (pendukung) & C3--C4 \\
\addlinespace
\CPMK{4} & Mengintegrasikan aset audio dan visual ke dalam interaksi
\istilah{game}. &
\CPL{04} (utama), \CPL{03} (pendukung) & C3--C5 \\
\bottomrule
\end{tabularx}
\end{table}

\begin{table}[htbp]
\centering
\caption{Matriks hubungan CPL dan CPMK}
\label{tab:matriks-cpl-cpmk-buku}
\begin{tabularx}{\textwidth}{@{}L{4.4cm}@{\extracolsep{\fill}}C{2.3cm} C{2.3cm} C{2.3cm} C{2.3cm}@{}}
\toprule
\headrow & \textbf{\CPMK{1}} & \textbf{\CPMK{2}} & \textbf{\CPMK{3}} & \textbf{\CPMK{4}} \\
\headrow & \footnotesize \istilah{Game loop} \& fisika & \footnotesize AI NPC & \footnotesize \istilah{Engine} \& level & \footnotesize Aset audio-visual \\
\midrule
\CPL{03} --- Algoritma dan bahasa & \Utama & \Utama & \Pendukung & \Pendukung \\
\CPL{04} --- Desain, implementasi, evaluasi & \Pendukung & \Pendukung & \Utama & \Utama \\
\bottomrule
\end{tabularx}
\tabsource{Keterangan: \textbf{Utama} = CPMK merupakan perwujudan langsung CPL;
Pendukung = CPMK memperkuat pencapaian CPL.}
\end{table}

\clearpage

\section*{Sub-CPMK}

Setiap CPMK dijabarkan menjadi Sub-CPMK yang lebih spesifik dan terukur.
Setiap Sub-CPMK pada tabel berikut menopang satu bab inti pada buku ini.

\begin{table}[htbp]
\centering
{\setlength{\tabcolsep}{4pt}\footnotesize%
\caption{Sub-CPMK, tingkat taksonomi Bloom, dan CPMK induknya}
\label{tab:subcpmk-buku}
\begin{tabularx}{\textwidth}{@{}C{1.5cm} Y C{1.6cm} C{1.5cm} C{1.7cm}@{}}
\toprule
\headrow \textbf{Kode} & \textbf{Rumusan Sub-CPMK} & \textbf{Bloom} & \textbf{CPMK} & \textbf{CPL} \\
\midrule
1.1 & Menjelaskan struktur \istilah{game loop} (input, \istilah{update},
\istilah{render}) dan peran manajemen waktu dalam permainan. &
C2 & \CPMK{1} & \CPL{03} \\
\addlinespace
1.2 & Merancang dan mengimplementasikan gerak objek berbasis \istilah{delta
time} dan vektor pada Pygame. & C3 & \CPMK{1} & \CPL{03} \\
\addlinespace
1.3 & Menerapkan logika fisika sederhana (gravitasi, lompatan, pantulan)
dalam simulasi 2D. & C3 & \CPMK{1} & \CPL{03} \\
\addlinespace
1.4 & Mendeteksi dan menangani \istilah{collision} menggunakan \code{Rect}
dan/atau \istilah{mask} pada \istilah{sprite}. &
C4 & \CPMK{1} & \CPL{03}, \CPL{04} \\
\addlinespace
2.1 & Merancang perilaku NPC sederhana (misalnya \istilah{chase})
berdasarkan posisi pemain. & C3 & \CPMK{2} & \CPL{03} \\
\addlinespace
2.2 & Mengimplementasikan \istilah{Finite State Machine} (idle, patrol,
chase, attack/flee) untuk NPC. & C3 & \CPMK{2} & \CPL{03} \\
\addlinespace
2.3 & Menerapkan \istilah{pathfinding} berbasis grid dengan algoritma A*
untuk navigasi NPC. & C4 & \CPMK{2} & \CPL{03} \\
\addlinespace
2.4 & Mengevaluasi perilaku NPC hasil kombinasi FSM dan A* terhadap
skenario permainan. & C5 & \CPMK{2} & \CPL{03}, \CPL{04} \\
\addlinespace
3.1 & Menavigasi antarmuka Unity dan menyusun \istilah{project},
\istilah{scene}, serta \istilah{GameObject} dasar. &
C3 & \CPMK{3} & \CPL{04} \\
\addlinespace
3.2 & Membangun level 2D/3D sederhana dengan \istilah{collider},
\istilah{Rigidbody}, dan kontrol pemain. & C4 & \CPMK{3} & \CPL{04} \\
\addlinespace
4.1 & Mengimpor dan mengonfigurasi aset visual (\istilah{sprite}/model,
animasi, UI) ke dalam \istilah{scene}. & C3 & \CPMK{4} & \CPL{04} \\
\addlinespace
4.2 & Mengintegrasikan aset audio dan memicu interaksi audio-visual
berdasarkan peristiwa \istilah{game}. & C4 & \CPMK{4} & \CPL{04} \\
\addlinespace
4.3 & Mengevaluasi prototipe \istilah{game} melalui \istilah{playtesting}
dan memperbaiki integrasi level serta aset. & C5 & \CPMK{4} & \CPL{04} \\
\bottomrule
\end{tabularx}}
\end{table}

\clearpage

\section*{Peta Bab, Materi, dan Minggu Perkuliahan}

Tabel berikut menghubungkan setiap bab pada buku ini dengan materi, Sub-CPMK,
dan minggu pelaksanaan pada RPS. Kolom \istilah{Teknologi} menunjukkan
perkakas utama yang dipakai pada bab tersebut.

\begin{table}[htbp]
\centering
{\setlength{\tabcolsep}{4pt}\footnotesize%
\caption{Peta bab buku ajar terhadap materi, Sub-CPMK, dan minggu RPS}
\label{tab:peta-bab}
\begin{tabularx}{\textwidth}{@{}C{1.2cm} L{4.3cm} C{1.5cm} C{1.8cm} C{1.5cm} Y@{}}
\toprule
\headrow \textbf{Bab} & \textbf{Judul Bab} & \textbf{Materi} & \textbf{Sub-CPMK} & \textbf{Minggu} & \textbf{Teknologi} \\
\midrule
I   & Pendahuluan dan Orientasi Buku & orientasi & --- & 1 & --- \\
II  & Landasan Teori dan Konsep Dasar & fondasi & --- & 1 & --- \\
III & Pengantar Pemrograman Game dan Persiapan Lingkungan & 1 & menuju 1.1 & 1 & Python/Pygame \\
IV  & \istilah{Game Loop}, \istilah{Event Handling}, dan \istilah{Clock} & 2 & 1.1 & 2 & Python/Pygame \\
V   & Gerak Objek, \code{Vector2}, Input, dan \istilah{Delta Time} & 3 & 1.2 & 3 & Python/Pygame \\
VI  & Fisika Sederhana: Gravitasi, Lompatan, dan Pantulan & 4 & 1.3 & 4 & Python/Pygame \\
VII & \istilah{Sprite}, \istilah{Group}, dan \istilah{Collision Detection} & 5 & 1.4 & 5 & Python/Pygame \\
VIII& Pengantar AI NPC dan Perilaku \istilah{Chase} & 6 & 2.1 & 6 & Python/Pygame \\
IX  & \istilah{Finite State Machine} untuk NPC & 7 & 2.2 & 7 & Python/Pygame \\
X   & \istilah{Pathfinding} Grid: dari BFS/Dijkstra ke A* & 8 & 2.3 & 8 & Python/Pygame \\
XI  & Integrasi FSM dan A* (Mini Proyek AI) & 9 & 2.4 & 9 & Python/Pygame \\
---  & \itshape Ujian Tengah Semester & --- & 1.1--2.4 & 10 & --- \\
XII & Pengenalan Unity: \istilah{Project}, \istilah{Scene}, \istilah{GameObject} & 10 & 3.1 & 11 & Unity \\
XIII& Membangun Level, \istilah{Physics}, dan Kontrol Pemain & 11 & 3.2 & 12 & Unity \\
XIV & Impor Aset Visual, Animasi, dan UI Dasar & 12 & 4.1 & 13 & Unity \\
XV  & Integrasi Audio dan \istilah{Trigger} Interaksi & 13 & 4.2 & 14 & Unity \\
XVI & Evaluasi Prototipe, \istilah{Playtesting}, dan Presentasi & 14 & 4.3 & 15 & Unity \\
---  & \itshape Ujian Akhir Semester & --- & 3.1--4.3 & 16 & --- \\
XVII& Evaluasi, Refleksi, dan Integrasi Kompetensi & penutup & 1.1--4.3 & --- & Python \& Unity \\
\bottomrule
\end{tabularx}}
\tabsource{Sumber: disusun dari RPS mata kuliah \NamaMK\ (\KodeMK).}
\end{table}

\begin{catatabox}
Bagian akhir buku memuat empat lampiran pendukung: Lampiran~A rubrik
penilaian lengkap, Lampiran~B contoh tugas dan lembar kerja penilaian,
Lampiran~C glosarium istilah, dan Lampiran~D daftar sumber belajar
tambahan. Peta bobot penilaian terhadap CPMK diuraikan pada Lampiran~A.
\end{catatabox}

\clearpage
